% This chapter describes the methodology, tools, and implementation choices
% used during the development of the AIS anomaly visualisation prototype.

\textit{This chapter outlines the development process, tools and frameworks used, system architecture, and testing strategies employed throughout the project.}

%=====================================================================
% 3.1  Development Methodology
%=====================================================================
\section{Development Methodology}

% FOCUS: Describe HOW you organised the project. Sprint structure,
% meeting routines, roles, and task tracking.
% Mirror the chess thesis Section 3.3 for level of detail.

%The team used an Agile approach with the Scrum framework to organise the project. Work was divided into sprints of one to two weeks, where sprint length was adjusted depending on workload from other courses.
% Unsure about what Arne should be called. Since he controls everything so just put him down as supervising lecturer.
\subsection{Scrum Framework}
The Scrum framework, described in Section 2.X, was adopted to structure the development process. The framework was introduced by the supervising lecturer, who provided pre-filled sprint templates and documentation to support its application for groups. This lowered the time investment of adopting the system and methodology, while ensuring that sprint artefacts such as the backlog and retrospectives were maintained throughout the project. Scrum was also well suited given the projects biweekly meetings with the supervisor and the customer at Kystverket. Each sprint would ensure that further development in the project could be presented and discussed, allowing scope adjustments to be made based based on the feedback. Should technical issues occur, such as limitations or issues with the AIS dataset or uncertainty around feasibility or scope for specific visualisation features. Then the iterative structure allows focus on the backlog until discussions about scope or features are had with the customer at Kystverket rather than following a fixed plan.

\subsection{Daily Stand-Ups}
Daily stand-up meetings were established as a core part of the sprint work flow, in line with the Scrum practices described in Section 2.X. The primary purpose of the stand-ups is to maintain accountability between team members and to further facilitate collaborative problem-solving between team members depending on problems blocking further development. In a two-person team, the risk of one member being stuck on a problem without the other member being aware is likely. Therefore by utilizing a daily meeting it allows such issues to be discussed early and not let it stop further project development. They can also be utilized as a daily forced break for the team. Sustained work without regular breaks is not equated to consistent productivity, while also being a way to quickly discuss a problem for further collaborative work after the stand-up meeting.


%Every day the team had a stand-up meeting at 12:00 where members discussed issues they had run into, what they were working on, and what to focus on next. However a lot of days we had multiple daily stand ups to clear our head and keep motivation high.

The project spanned from January 2026, with delivery on the 20th of May. 

The project was initially started by a team of three members. However one member withdrew from the project early in the semester due to personal reasons. The remaining two members, continued development and divided the responsibility between them. One member served as Scrum Master and maintained the sprint documentation, while the other handled meeting notes. All other responsibilities including development, testing, and design were shared equally.
% TODO: Add about issues with workload and how we handled it.

Biweekly meetings were held with the supervisor to review academic progress and get guidance. We also mostly had biweekly meetings with the customer at Kystverket which were arranged to present progress, get feedback, and clarify scope.
% TODO: Maybe talk about how they were busy some times also wondering how to word kystverket.

\subsection{Task Management}

Task tracking was handled using GitHub Projects, which provided a shared Kanban-style board across the team's three repositories: one for the frontend, one for the backend, and one for administrative documents like meeting notes and sprint documentation. Issues were created and tracked through workflow stages to maintain workload transparency and motivation.
% TODO: Maybe add link and theory to theory chapter
% https://docs.github.com/en/issues/planning-and-tracking-with-projects/customizing-views-in-your-project/customizing-the-board-layout

%=====================================================================
% 3.2  Tools and Platforms
%=====================================================================
\section{Tools and Platforms}

\subsection{Development Tools}
\begin{itemize}
\item \textbf{Visual Studio Code:} Primary development environment used by all team members.
\item \textbf{Git \& GitHub:} Used for version control, branching, pull requests, and task management via GitHub Projects.
\item \textbf{Docker:} The backend was containerised using Docker to ensure a consistent and reproducible runtime environment.
\item \textbf{Swagger:} Used for API documentation and testing of backend endpoints.
\end{itemize}

\subsection{Design Tools}
\begin{itemize}
\item \textbf{Figma:} Used for designing wireframes and UI mockups for the web application.
\end{itemize}

\subsection{Communication Tools}
\begin{itemize}
\item \textbf{Discord:} Primary communication platform for team discussions, quick coordination, and informal meetings.
\item \textbf{Slack:} Used for communication with the customer (Kystverket), including asking technical questions and discussing issues.
\item \textbf{Outlook:} Used for planning meetings with Supervisor.
\end{itemize}

\subsection{AI Tools}
\begin{itemize}
\item \textbf{ChatGPT \& Gemini:} Primarily used for technical issues, understanding technologies, and spell checking. 
\end{itemize}

% TODO: Talk with other group member about their ai usage to make sure nothing is missing
% TODO: Add other tools


%=====================================================================
% 3.3  Technology Stack
%=====================================================================
\section{Technology Stack}

The technology stack was primarily decided by the customer's existing infrastructure. Since the prototype is intended to work with Kystverket's systems, using their technology choices ensures compatibility. It also allows the team to leverage the customer's expertise if questions arise, since this is a prototype using their techstack.

\subsection{Frontend}
\begin{itemize}
\item \textbf{React \& TypeScript:} Core frontend framework and language, chosen as a requirement from the customer.
\item \textbf{MapLibre GL JS:} Mapping engine used for rendering interactive, high-performance maps with WebGL acceleration.
%\item \textbf{Vite:} 
% TODO: Go over as a group to make sure no libraries or dependencies are missing.

\end{itemize}

\subsection{Backend}
\begin{itemize}
\item \textbf{PostgreSQL:} The database storing the AIS anomaly test data, provided by the customer.
\item \textbf{Swagger:} API documentation is auto-generated using Swagger annotations in the Go source code. A Swagger UI endpoint is available for exploring and testing the API.
\item \textbf{Docker:} Used for containerising the backend services.
% TODO: Mention other things used in the backend making sure nothing is missing.
\end{itemize}

The backend was originally provided by Kystverket as a working codebase pre-populated with test data, since actual AIS anomaly data is confidential. The team forked this repository and updated it throughout the project. The most notable addition was a master query endpoint described in Section 3.4.1
% TODO: Describe the test dataset briefly (rough size, time range, types of anomalies if known).

%=====================================================================
% 3.4  System Architecture
%=====================================================================
\section{System Architecture}

% FOCUS: Describe how the pieces fit together. Include a diagram if possible.
% The chess thesis has sequence diagrams, class diagrams, package diagrams —
% you should have at least one architecture diagram here.

% TODO: Add a figure showing the overall system architecture.
% A simple box diagram: Database <-> Backend (Go/Docker) <-> REST API <-> Frontend (React/MapLibre)

The system follows a client-server architecture. The backend is written in Go and runs inside a Docker container that exposes a REST API. Swagger is used for API documentation. The frontend is built with React and TypeScript and retrieves the data it needs through API calls to the backend.

\subsection{Backend and Data Flow}

The backend infrastructure and test database were provided by the customer, pre-populated with test data since actual AIS anomaly data is confidential. The team's main contribution to the backend was a master query endpoint. This endpoint accepts all filter parameters (anomaly type, mmsi, and time range.) and returns only the matching data. Instead of getting all the anomaly groups and filtering on the frontend, all variables get sent to the backend which queries the data again. This was done since keeping old data and filtering client-side was slow and inefficient.

% TODO: Maybe include how data is returned and used? Geojson and such

\subsection{Frontend Architecture}
The frontend architecture follows a client-driven data flow where data is retrieved from the backend through REST API calls using Axios. The received data is stored in React state as a geoJSON featureCollection, which is passed to individual Deck.gl layers responsible for rendering the visualisations. This is done after a map instance is created through a map container class, using maplibre. deck.gl is subsequently used for rendering data points overlaid visually on the map. Interactivity is handled through React components, where state changes control filtering and updating of the rendered layers. Individual Deck.gl layers also implement click event handlers, which update application state and trigger the rendering of additional information components or different visualization layers. 

% TODO: Describe the React component structure at a high level.
% - How is routing handled?
% - What are the main pages/views? (Map view, filter panel, detail cards, etc.)
% - How does state management work for filters and map state?

%=====================================================================
% 3.5  Visualisation Implementation
%=====================================================================
\section{Visualisation Implementation}

% FOCUS: This is your core contribution — give it proper space.
% Describe HOW you implemented each visualisation technique from your Theory chapter.

\subsection{Map Setup and Rendering}

The base map is implemented using MapLibre GL JS, as the engine for rendering the map. based on a vector tile set from openFreeMap  \url{https://tiles.openfreemap.org}. \label{sec:vector tiles} The tile set chosen was bright in color. Deck.gl was chosen for rendering data points on the map, with points being overlaid on top of the base map



% TODO: Describe the base map setup:
% - What vector tile source did you use? (customer-provided? OpenMapTiles? Maptiler?)
% - Any custom map styles?
% - Initial map center, zoom level, bounds?

\subsection{Clustering}


As mentioned in Section \ref{sec:supercluster}, Supercluster.js is used for calculating clusters based on zoom level and the current map bounds. Supercluster calculates clusters based on the base dataset acquired from the API, calculating all clusters on application load. Map interactions such as zooming and dragging result in clusters being retrieved for the current zoom level and map bounds. The clustering data itself is used with deck.gl,Section \ref{sec:deckgl}. Visualization through deck.gl is accomplished through creating a custom layer, consisting of a label-layer, whose purpose is to be a point counter for the amount of anomalies inside the cluster, this is put over, a scatterplot-layer. section \ref{sec:scattermap}, which combines to create the clustering layer.This layer is re-rendered whenever the clustering output changes due to updates in zoom level or map bounds. Furthermore does clustering visualization have dynamics sizing based on amount of anomalies inside a cluster, leading to clusters being scaled proportional to their anomaly count. 

% TODO: Add technical details:
% - What clustering radius/parameters did you use?
% - How are cluster colours determined (step expressions? interpolation?)
% - Any custom cluster styling?
% - add figure

\subsection{Heat Map}
A heat map layer was implemented for visualization of the full dataset. \ref{sec:heatmap}. deck.gl was used for heat map visualization, where the library provided the methods to create the heat map, with the data provided being the base API data of the full dataset. 
% TODO: Describe the MapLibre heat map layer configuration:
% - What colour ramp did you use?
% - How does the weight/intensity scale with zoom?
% - Any specific property used for weighting (e.g. anomaly severity)?


\subsection{Anomaly Group Layer}
The anomaly group layer, is a visualization layer that gives a detailed representation of a single anomaly group, section \ref{sec:anomalies and anomaly groups}. The anomaly group visualization layer is created through a custom deck.gl layer. This layer consists of a signal strength layer, using data from the individual anomalies inside the group, that are connected to a satellite or base station with a signal strength number value, normalized between 0 and 1. The signal strength layer adds lines from the source anomaly, to the receiving base station and or satellite, indicating the strength of received AIS-message.   





\subsection{Frontend Components}


\subsection{Scope and Feature Prioritisation}

The project scope was decided through ongoing discussions with the customer. Hexagonal binning was initially identified as a desired feature in the project brief. However due to time constraints it was not implemented in the final prototype. The team prioritised clustering, heat maps, and the filter system based on customer feedback since these provided the most immediate analytical value. 
% TODO: Include what we cut / didnt focus on
% NOTE TO SELF: Your Introduction (Section 1.2 Objective) lists hexagonal binning
% and advanced temporal analysis as goals. Since hex binning was not implemented
% and temporal analysis is limited to time-range filtering, you should either:
%   (a) soften the Introduction wording (e.g. "investigate" instead of "implement"), or
%   (b) address the gap explicitly in Discussion (recommended — explain the
%       scoping decision made with the customer).


\subsection{Filter System}



\subsection{Temporal Filtering}



\subsection{Interactive Features}


%=====================================================================
% 3.6  Design Process
%=====================================================================
\section{Design Process}

% FOCUS: Wireframes, UI/UX decisions, user feedback.

\subsection{Wireframes and UI Design}


\subsection{User Feedback}


%=====================================================================
% 3.7  Testing
%=====================================================================
\section{Testing}

% FOCUS: What you tested, how, and why certain approaches were chosen.

\subsection{Backend API Testing}


\subsection{Frontend Testing}


\subsection{Manual and Performance Testing}

%=====================================================================
% End of Methods chapter
%=====================================================================
